\section{Motivation}

\subsection{Digitizing Vietnamese Documents}

Vietnam has undergone rapid administrative modernization over the past two decades, driven by government-led digital transformation initiatives that aim to replace manual, paper-based workflows with electronic document management systems across public and private sectors \cite{Decision749QDTTg_2020}. This transition affects an extraordinarily broad range of document types: tax invoices, procurement records, official correspondence, land registration certificates, health records, educational transcripts, court decisions, and legislative texts, all of which are currently stored in paper form in archives accumulated over many decades of governance \cite{MST_Egov_Strategy_2021}. The sheer volume of existing paper documentation, combined with the ongoing generation of new paper records in agencies that have not yet completed their digital transitions, creates a pressing demand for automated document digitization at scale \cite{Kaiser2022}.

OCR serves as the key enabling technology for this digitization effort. Without OCR, scanned documents remain image files that cannot be searched, indexed, or used as inputs to downstream information extraction systems. With OCR, scanned documents become machine-readable text that can be ingested into databases, processed by information retrieval pipelines, and made accessible through full-text search interfaces. The economic and administrative value of this transformation is substantial: automated OCR processing can reduce manual data entry costs by orders of magnitude, accelerate document retrieval times from hours to milliseconds, and enable the construction of queryable national archives that would otherwise require prohibitive amounts of human labor to maintain \cite{Palm2017}.

The digitization of Vietnamese historical and archival documents carries significant cultural and scholarly value. Libraries, museums, and academic institutions hold large collections of historical Vietnamese texts, newspaper archives, colonial-era administrative records, and literary manuscripts that are deteriorating due to age and environmental factors. Digitization through OCR preserves these documents in a form that can be archived, replicated, and made accessible to researchers and the public worldwide. The linguistic and historical content of such archives is irreplaceable once lost, giving OCR development for Vietnamese a preservation dimension that extends beyond immediate administrative utility.

The educational sector similarly benefits from large-scale OCR capability. Textbooks, examination records, student transcripts, and pedagogical materials exist predominantly in printed form across Vietnamese educational institutions, and their digitization enables new forms of accessibility, translation, and adaptive learning that are impossible with image-only storage. The combination of administrative, archival, and educational demand establishes a compelling and multi-dimensional motivation for investing in high-quality Vietnamese OCR technology.

\subsection{Challenges in Vietnamese OCR}

Vietnamese presents a unique set of challenges for optical character recognition that extend well beyond those encountered in most Latin-script languages and that are not adequately addressed by general-purpose OCR systems developed without Vietnamese-specific adaptation \cite{Le2025}. These challenges arise from two interconnected sources: the exceptional orthographic complexity of the Vietnamese writing system, and the limited availability of annotated data resources for training and evaluating recognition models in this language \cite{Le2025}.

At the orthographic level, Vietnamese employs a diacritical system of unusual density and lexical significance. Unlike most Latin-script languages where diacritical marks play a limited and relatively predictable role, Vietnamese superimposes two independent layers of diacritical modification onto its base characters simultaneously. The first layer consists of vowel quality markers that alter the phonetic identity of the base vowel, producing three distinct characters from the base vowel \textit{a} (plain \textit{a}, circumflex \textit{a}, and breve \textit{a}), with analogous modifications applying to several other base vowels and yielding seventeen distinct vowel surface forms in total \cite{Nguyen2012}. The second layer consists of tonal markers encoding one of six lexically contrastive tones through five diacritical marks and the absence of any mark for the level tone. Because these two layers of modification can co-occur on a single character, the resulting Vietnamese character inventory is several times larger than that of unaccented Latin-script languages, posing a direct challenge to recognition models trained on standard Latin benchmarks \cite{Nguyen2021}.

From the perspective of OCR model design, this orthographic complexity produces two primary failure modes. The first is character-level confusion, in which visually similar composite glyphs are misclassified due to the small spatial extent of distinguishing diacritical strokes relative to the base character body. At low image resolutions or under degraded scan conditions, the difference between a dot below and a hook above, or between a circumflex and a breve, can be reduced to a handful of pixels, falling below the discriminative capacity of feature extractors that were not specifically trained on Vietnamese text \cite{Nguyen2021}. The second failure mode is segmentation ambiguity, in which the spatial extent of a composite Vietnamese character overlaps with neighboring characters or with the ascender and descender zones of adjacent text lines, causing detection-based systems to either merge distinct characters or fragment a single character across multiple detected regions \cite{Nguyen2021}.

Beyond orthographic complexity, the diversity of Vietnamese document formats introduces additional technical challenges for both detection and recognition stages. Vietnamese administrative documents span a wide range of layouts, including single-column official correspondence, multi-column legal texts, tabular financial records, receipts with characteristic formatting structures, and identity documents with fixed field templates. Real-world documents may be captured through flatbed scanning, camera photography, or mobile phone imaging, each introducing different geometric distortions, illumination conditions, and resolution characteristics. Text lines in scanned documents are frequently not perfectly horizontal due to page skew introduced during scanning, and documents with physical curvature produce locally varying text orientations that challenge axis-aligned detection systems. The diversity of Vietnamese font styles across formal administrative typefaces, commercial display fonts, and historical typographic conventions further expands the distribution that a robust OCR system must cover.

A compounding challenge is the limited availability of annotated Vietnamese OCR datasets suitable for training modern deep learning systems. Compared to the rich benchmark ecosystems surrounding English and Chinese OCR, Vietnamese-specific datasets remain sparse in terms of scale, annotation granularity, and document type diversity \cite{Nguyen2021}. The most linguistically diverse publicly available Vietnamese OCR benchmark, VinText \cite{Nguyen2021}, comprises approximately 2,000 real images, which is several orders of magnitude smaller than the large-scale multilingual datasets used to train state-of-the-art multilingual OCR systems such as MLT \cite{Nayef2017}. The absence of character-level annotations in most Vietnamese datasets further restricts the ability to train and evaluate detection components at fine granularity, and the lack of standardized evaluation protocols across existing datasets makes cross-study comparison unreliable.

The scarcity of annotated data is directly exacerbated by the orthographic complexity of Vietnamese: annotating Vietnamese text requires annotators who can reliably distinguish between visually similar tone-diacritic combinations, which demands higher annotator skill and more careful quality control than annotation of unaccented Latin text. This creates a self-reinforcing bottleneck in which the language most in need of specialized OCR systems is also the most difficult to annotate efficiently, limiting the scale of training data that can be produced for a given annotation budget.

\subsection{The Importance of Accurate Text Recognition in Vietnamese}

The importance of recognition accuracy in Vietnamese OCR systems is not merely a matter of benchmark performance but is directly tied to the practical consequences of recognition errors in the deployment contexts where these systems are most needed. Vietnamese is a tonal language in which diacritical marks carry lexical meaning: altering or omitting a single diacritical mark on a character does not produce a misspelling of the intended word but an entirely different word with a potentially unrelated meaning \cite{Le2025}. For example, the syllables \textit{ma}, \textit{mà}, \textit{má}, \textit{mả}, \textit{mã}, and \textit{mạ} (differing only in tonal mark) denote distinct words corresponding to ghost, but, cheek, tomb, horse and rice seedling respectively \cite{Nguyen2019}. A recognition error that substitutes one tonal mark for another therefore does not corrupt a character; it changes the identity of the word, and the corrupted output may be another valid Vietnamese word that passes any lexical validity check without revealing the underlying error \cite{Pham2017}.

This property of Vietnamese orthography makes recognition errors both more consequential and harder to detect than in most other Latin-script languages. In the context of administrative document processing, recognition errors on invoice fields, contract terms, official names, or legal identifiers can introduce downstream errors that propagate through automated extraction pipelines without triggering any alert. The consequences range from financial discrepancies in automated invoice processing to legal ambiguities in digitized contract records and identity mismatches in administrative databases \cite{Park2019}. Vietnamese government agencies processing tax invoices, procurement records, and official correspondence through OCR-based digitization pipelines are particularly exposed to these risks, as the orthographic complexity of Vietnamese text imposes additional demands on the recognition component and recognition errors on diacritical marks can produce semantically distinct output that corrupts downstream information extraction \cite{Huang2019}.

In historical archive digitization, the stakes are different but equally high. OCR errors in the transcription of historical documents create persistent inaccuracies that, once introduced into searchable archives, are difficult to identify and correct without reviewing the original scanned images. A digitized historical corpus with significant recognition error rates is not merely imperfect but may be actively misleading, attributing incorrect words or names to historical figures and creating artifacts in computational analyses that rely on the corpus as primary data. The cultural and scholarly value of accurate historical transcription makes recognition accuracy a prerequisite for the usefulness of any archival OCR system, not simply a performance optimization.

In information retrieval systems, recognition accuracy directly determines recall, the fraction of relevant documents retrieved in response to a query. If the OCR output for a document contains character errors, keyword searches for terms that appear in that document may fail to retrieve it, effectively making the document invisible to the information retrieval system. Given that Vietnamese administrative archives contain documents that may be retrieved for legal, historical, or policy purposes years or decades after their creation, the long-term consequences of recognition errors that suppress relevant retrieval are substantial. High recognition accuracy is therefore not a quality-of-life improvement but a functional requirement for OCR-based information retrieval to serve its intended purpose.

These considerations collectively establish that Vietnamese OCR is a problem of practical importance that demands solutions commensurate with the difficulty of the language and the stakes of the deployment context. The remainder of this thesis is dedicated to developing and evaluating such a solution.

\subsection{Recent Applications in Document OCR}

The practical relevance of optical character recognition extends well beyond academic benchmarks, with deployed OCR systems now underpinning a wide range of real-world document processing workflows across industries and public institutions. The maturation of deep learning-based detection and recognition architectures has enabled OCR to move from controlled laboratory settings into operational environments characterized by document diversity, scale requirements, and strict accuracy demands. This section surveys three representative application domains that illustrate the breadth of contemporary OCR deployment: administrative and financial document processing, identity document recognition, and historical archive transcription. Each domain presents a distinct set of technical challenges that contextualize the practical value of the Vietnamese OCR system developed in this thesis.

\subsubsection{Administrative and Financial Document Processing}
The digitization of administrative and financial documents, including invoices, receipts, purchase orders, and tax forms, represents one of the most commercially mature applications of OCR technology. Organizations processing large volumes of such documents have strong economic incentives to automate extraction of structured information such as vendor names, dates, line items, and totals, reducing reliance on manual data entry and accelerating downstream business processes. Early OCR-based approaches to this problem relied on template matching against fixed document layouts, but the diversity of invoice formats across vendors and jurisdictions made template-based systems brittle and expensive to maintain \cite{Palm2017}.

More recent work has framed invoice and receipt understanding as a joint OCR and information extraction problem, in which text detection and recognition are coupled with named entity recognition and document layout analysis to produce structured output from unstructured document images. The CORD dataset \cite{Park2019} and the SROIE benchmark \cite{Huang2019} established standardized evaluation frameworks for receipt OCR and key information extraction, driving the development of end-to-end trainable systems that simultaneously localize text, transcribe it, and classify it into semantic categories. LayoutLM \cite{Xu2020} extended transformer-based language models to the document understanding domain by incorporating two-dimensional positional embeddings that encode the spatial coordinates of each text token alongside its semantic embedding, enabling the model to exploit document layout structure as an additional signal for information extraction. Subsequent iterations including LayoutLMv2 \cite{Xu2022} and LayoutLMv3 \cite{Huang2022} further integrated visual features from document images directly into the pretraining objective, producing unified multimodal representations that have become the dominant paradigm for structured document understanding

The relevance of this application domain to Vietnamese OCR is direct. Vietnamese government agencies and enterprises generate large volumes of administrative documents including tax invoices, procurement records, and official correspondence, all of which are candidates for automated digitization. The orthographic complexity of Vietnamese text imposes additional demands on the OCR component of such pipelines, as recognition errors on diacritical marks can produce semantically distinct output that corrupts downstream information extraction.

\subsubsection{Identity Document Recognition}
Identity document recognition, encompassing the automated extraction of personal information from national identity cards, passports, driving licenses, and similar credentials, constitutes another high-value OCR application domain with stringent accuracy and reliability requirements. Errors in recognized identity fields can have serious consequences in contexts such as border control, financial onboarding, and voter registration, making this one of the most demanding real-world OCR deployment scenarios.

Early systems exploited the structured layouts of standardized identity documents, using template registration to locate predefined field regions before applying OCR within each. Machine-readable zones (MRZ) provided a constrained recognition target well-suited to classical techniques, given their limited character set and fixed typography. However, the diversity of identity document formats across countries and the presence of security printing features, holographic overlays, and variable background textures have motivated the adoption of deep learning-based approaches that generalize across document types without requiring per-format template engineering. The MIDV-500 dataset \cite{Bulatov2019} established a standardized benchmark for identity document analysis by providing annotated samples across 50 document types from different countries under varied capture conditions including flat, curved, and perspective-distorted document images, enabling more rigorous evaluation of generalizable identity document recognition systems.

More recent work has framed identity document recognition as an end-to-end structured information extraction problem, combining text detection, recognition, and field classification within a unified pipeline \cite{Tran2019}. Rather than treating each processing stage independently, these systems jointly optimize detection and recognition objectives, enabling the model to leverage layout structure as a supervisory signal for improving transcription accuracy across variable document formats. For Vietnamese identity documents specifically, the transition from the old nine-digit citizen identification card to the new twelve-digit CCCD format introduced new layout conventions and font specifications that required OCR systems to be retrained or fine-tuned, highlighting the importance of adaptable recognition architectures \cite{Nguyen2022}. The cascaded detection-recognition design adopted in this thesis is well suited to this setting, as the modular architecture allows the detection and recognition components to be independently fine-tuned when document specifications change.

Across both application domains surveyed in this section, a consistent pattern emerges: the performance of deployed OCR systems is ultimately bounded by the quality of the underlying text detection and recognition components, and improvements to these core components translate directly into downstream application quality. This observation reinforces the motivation for the thesis at hand, which targets the detection and recognition pipeline as the primary locus of improvement for Vietnamese document OCR.

\subsection{Vietnamese-Specific OCR Challenges and Solutions}
Vietnamese presents a unique set of challenges for optical character recognition that extend well beyond those encountered in most Latin-script languages, and that are not adequately addressed by OCR systems designed for general multilingual use. While the deep learning architectures surveyed in the preceding section have demonstrated strong performance across a broad range of scripts and document types, their direct application to Vietnamese text without domain-specific adaptation consistently yields suboptimal results. These shortcomings stem from two interconnected factors that distinguish Vietnamese OCR from the general case.

The first factor is linguistic and orthographic in nature. Vietnamese employs a diacritical system of exceptional complexity, in which multiple combining marks can co-occur on a single character to encode both vowel quality and lexical tone simultaneously. This orthographic density produces a character inventory significantly larger than that of unaccented Latin-script languages, and introduces visual ambiguities that challenge both detection and recognition components of an OCR pipeline. The second factor is resource-related. Despite growing interest in Vietnamese document digitization, the landscape of publicly available annotated datasets and standardized benchmarks for Vietnamese OCR remains considerably less developed than the corresponding ecosystems for English, Chinese, or Arabic, constraining the training and rigorous evaluation of recognition models.

These two factors are closely interrelated: the orthographic complexity of Vietnamese makes the annotation process more demanding and error-prone, which in turn limits the scale and coverage of available datasets, which in turn impedes the development of models robust enough to handle the full range of Vietnamese diacritical variation. This section examines both dimensions in turn, first analyzing the specific diacritical and tonal characteristics of Vietnamese that create difficulties for OCR models, and then surveying the existing datasets and benchmarks that have been developed to support Vietnamese text recognition research, identifying the gaps that motivate the data preparation strategy adopted in this thesis.

\subsubsection{Diacritical Marks and Tonal Complexity}
\label{sec:diacritical_marks_and_tonal_complexity}
Vietnamese presents a distinctively challenging target for optical character recognition due to its orthographic system, which overlays multiple layers of diacritical modification onto a Latin-derived base script. Unlike most Latin-script languages, where diacritics serve a limited and relatively predictable role, Vietnamese employs diacritics at two independent levels simultaneously: vowel quality markers that alter the phonetic identity of the base character, and tonal markers that specify one of six lexically contrastive tones \cite{Nguyen2012}. These two layers of modification can co-occur on a single character, producing composite glyphs that are visually dense and spatially compact, creating significant difficulties for both text detection and recognition systems.

At the level of vowel modification, Vietnamese augments several base vowels with diacritic strokes that fundamentally change their phonetic identity. The vowel \textit{a}, for instance, yields three distinct characters: \textit{a}, \textit{a with circumflex} (â), and \textit{a with breve} (ă), each representing a different phoneme and therefore a different lexical unit. Similar modification patterns apply to the vowels \textit{e}, \textit{o}, and \textit{u}, yielding a total of seventeen distinct vowel surface forms from seven base vowels. When tonal markers are superimposed on these modified vowels, the resulting character set expands substantially. Vietnamese defines six tones encoded through five diacritical marks, namely the grave, acute, hook above, tilde, and dot below, plus the absence of any tone mark for the level tone. Combining all vowel modifications with all tonal possibilities yields a Vietnamese character inventory that is several times larger than that of unaccented Latin-script languages, posing a direct challenge to recognition models trained on standard Latin benchmarks \cite{Nguyen2021}.

From the perspective of OCR model design, this orthographic complexity manifests in two principal failure modes. The first is character-level confusion, where visually similar composite glyphs are misclassified as one another due to the small spatial extent of distinguishing diacritical strokes relative to the base character body. At low image resolutions or under degraded scan conditions, the difference between a dot below and a hook above, or between a circumflex and a breve, can be reduced to a handful of pixels, well below the discriminative capacity of feature extractors trained on higher-resolution data \cite{Nguyen2021}. The second failure mode is segmentation ambiguity, where the spatial extent of a composite Vietnamese character overlaps with neighboring characters or with the ascender and descender zones of adjacent text lines, causing segmentation-based detectors to either merge distinct characters or fragment a single character across multiple segments \cite{Nguyen2021}. This segmentation problem is particularly acute in printed documents with tight inter-line spacing and diverse font styles, which fall squarely within the scope of this thesis.

Existing approaches to mitigating these challenges fall into two broad categories. The first encompasses character set redesign strategies, in which the model output vocabulary is explicitly constructed to include all valid Vietnamese composite characters as atomic units, rather than decomposing them into base characters and combining diacritics \cite{Nguyen2021}. This atomic treatment forces the recognition model to learn a direct mapping from visual glyph to full Unicode character, avoiding the error accumulation that arises when base character and diacritic predictions are made independently. The second category leverages linguistic constraints through language model integration, exploiting the fact that not all combinations of base character and diacritical marks are phonologically valid in Vietnamese, and that valid tone-vowel combinations follow predictable phonotactic rules that can be encoded as decoding constraints \cite{Nguyen2021}. The methodology adopted in this thesis draws on both categories, as detailed in Chapter \ref{chapter:methodology_chapter3}.

\subsubsection{Existing Vietnamese OCR Datasets and Benchmark}
\label{sec:existing_vnese_ocr_datasets_and_benchmark}

The availability of high-quality annotated datasets is a prerequisite for training and evaluating robust Vietnamese OCR systems, yet the Vietnamese OCR research community has historically been constrained by a limited and fragmented landscape of publicly available corpora. Compared to the rich benchmark ecosystems surrounding English and Chinese OCR, Vietnamese-specific datasets remain sparse in terms of both scale and diversity, a gap that has directly impeded the development of generalizable recognition models \cite{Nguyen2021}.

Among the earliest publicly available resources is the VOHTR2018 dataset \cite{Nguyen2018}, which targets offline handwritten Vietnamese text and provides word-level and line-level images annotated with ground truth transcriptions. While valuable for handwriting recognition research, VOHTR2018 does not address the printed document domain and its annotation granularity is insufficient for training text detection models that require bounding box supervision at the word or character level. Its annotations also do not generalize readily to scene text conditions, limiting its utility for the broader Vietnamese OCR task.

A more recent and comprehensive contribution is VinText \cite{Nguyen2021}, a Vietnamese scene text dataset comprising over 2,000 real images collected from diverse real-world environments including street signs, storefronts, menus, and printed documents. VinText provides word-level polygon annotations along with full transcriptions and constitutes the most linguistically diverse publicly available Vietnamese OCR benchmark to date. Its inclusion of natural scene images with complex backgrounds, arbitrary text orientations, and mixed font styles makes it substantially more challenging than synthetic corpora and more representative of deployment conditions. A key contribution of VinText is its dictionary-guided recognition framework, which explicitly incorporates Vietnamese lexical constraints to resolve diacritical ambiguities at inference time \cite{Nguyen2021}. However, its scale remains modest relative to large-scale multilingual benchmarks such as MLT \cite{Nayef2017} and LSVT \cite{Sun2019}, and its coverage of formal administrative document layouts is limited.

Across these datasets, several consistent gaps are apparent. First, no existing Vietnamese OCR benchmark provides sufficient coverage of the full range of document types encountered in administrative digitization workflows, including low-resolution scans, faded ink, and degraded print quality. Second, character-level annotation remains rare, with most datasets providing only word-level or line-level ground truth, which limits the ability to train and evaluate detection components at finer granularity. Third, the absence of a standardized evaluation protocol across datasets makes cross-paper comparison of reported results unreliable. These gaps collectively motivate the data preparation strategy adopted in this work, which combines synthetic generation with domain-specific augmentation to compensate for the scarcity of real annotated Vietnamese document images, as described in Chapter \ref{chapter:methodology_chapter3}.

\section{Proposed Approach: Deep Learning-Based \\ Vietnamese OCR}

\subsection{Approach}
\label{sec:approach_1_2_1}
The approach adopted in this thesis is grounded in the observation that the two fundamental subtasks of document OCR, spatial localization of text regions and sequential transcription of their contents, have distinct computational requirements that are best addressed by architectures specifically optimized for each. Rather than pursuing a single end-to-end model that attempts to solve both problems simultaneously within a shared representation, this work proposes a cascaded two-stage architecture in which a specialized detector and a specialized recognizer are independently trained and then connected through a well-defined geometric interface.

For the detection stage, this thesis adopts RT-DETR V2 \cite{Zhao2024}, a real-time transformer-based object detector that combines the spatial feature extraction efficiency of convolutional networks with the global context modeling capacity of transformer self-attention. RT-DETR is selected over convolutional alternatives such as EAST \cite{Zhou2017} and CRAFT \cite{Baek2019} for several reasons that are directly relevant to the Vietnamese document setting. Its anchor-free, set-prediction formulation eliminates the need to manually engineer aspect ratio priors for Vietnamese text lines, which vary considerably across document types and orientations. Its hybrid encoder captures both local character-level texture and global line-level structure within a single forward pass, without the receptive field limitations inherent to purely convolutional architectures. Furthermore, its strong out-of-the-box localization accuracy on general object detection benchmarks provides a favorable initialization for domain-specific fine-tuning on Vietnamese text data. To enable oriented bounding box prediction, which is necessary for tightly enclosing skewed and rotated text lines encountered in real Vietnamese documents, RT-DETR V2 is extended with an auxiliary angle regression head that produces sine-cosine angle predictions, coupled with a Gaussian Wasserstein Distance (GWD) loss \cite{Yang2022} that provides geometrically principled supervision over all five oriented box parameters jointly.

For the recognition stage, this thesis adopts PARSeq \cite{Bautista2022}, a Vision Transformer-based scene text recognizer that formulates character sequence prediction as a permutation language modeling problem. PARSeq is selected because its ViT encoder captures the fine-grained spatial detail required to distinguish Vietnamese diacritical marks that differ only by subtle stroke variations, while its permutation-trained decoder implicitly learns a bidirectional language model that can resolve character-level ambiguities using both preceding and following character context. This bidirectional reasoning capability is particularly valuable for Vietnamese, where the distinction between visually similar tone-diacritic combinations is frequently underdetermined by local visual evidence alone and requires contextual inference from the surrounding character sequence. PARSeq is adapted to the Vietnamese domain through a carefully designed vocabulary extension that treats each composite Vietnamese glyph as an atomic output token, class-weighted training that compensates for the severe frequency imbalance between common and rare diacritical combinations, and a phonological validity mask applied at decoding time to suppress linguistically impossible character outputs.

The combination of these two components into a cascaded pipeline represents the central technical contribution of this thesis. The modular design allows each component to be independently trained, evaluated, and improved without requiring full system retraining, which is a practical necessity given the limited availability of annotated Vietnamese OCR data. The interface between the two stages performs perspective-corrected cropping using the oriented bounding box predictions from the detector, converting arbitrarily rotated text line regions into upright rectangular crops before passing them to the recognizer. This cropping strategy is not merely an engineering convenience but is a direct operationalization of the theoretical argument that tight angular localization in the detection stage is a prerequisite for accurate diacritical recognition in the transcription stage.

\subsection{Pipeline Components and Workflow}

The end-to-end pipeline transforms a raw input document image into a structured collection of recognized text strings through five sequential processing stages, each implemented as an independent module with well-defined input and output interfaces.

The pipeline transforms a raw input document image into a structured collection of recognized text strings through five sequential stages: image ingestion and preprocessing, oriented text line detection, bounding box post-processing, perspective-corrected cropping, and character sequence recognition. Each stage is implemented as an independent module, enabling the system to be executed as a unified inference pipeline. The mechanics of each stage, including intermediate tensor shapes and inter-stage data contracts, are elaborated in Chapter~\ref{chapter:methodology_chapter3}.

\subsection{Model Optimization Strategy}

Achieving competitive performance on Vietnamese text with the proposed cascaded architecture requires targeted optimization of both pipeline components, since the general-purpose pretrained checkpoints for RT-DETR and PARSeq were not trained on Vietnamese document data and do not natively support the full Vietnamese composite character set. The optimization strategy applied to each stage is briefly summarized here; complete technical details are provided in Chapter 3.

For the detection stage, the primary optimization challenge is adapting a model pretrained on general object detection to the single-class task of Vietnamese text line localization with oriented bounding box output. This is addressed through a combination of domain-specific synthetic data generation and a carefully staged fine-tuning protocol. The synthetic data pipeline produces approximately 130,000 annotated document images spanning seven layout archetypes, including single-column administrative correspondence, multi-column texts, receipt formats, and fully rotated scatter layouts, with exact oriented bounding box annotations derived analytically from rendering parameters. A critical engineering contribution of this pipeline is the correction of an angle distribution collapse in the base generator, which ensures that the training set provides uniform angular coverage across the full $[0^\circ, 360^\circ)$ rotation range. The fine-tuning protocol employs a three-tier differential learning rate schedule that assigns distinct learning rates to the pretrained backbone, the transformer encoder-decoder, and the newly initialized angle regression head, combined with metric-driven progressive backbone unfreezing to protect pretrained representations during early training. The loss function combines standard detection losses with a cosine angle loss and the GWD loss weighted at $\lambda_{\text{gwd}} = 3.0$, reflecting the importance of angular supervision for downstream recognition quality.

For the recognition stage, the primary optimization challenges are vocabulary expansion and class imbalance. The standard PARSeq vocabulary of 94 printable ASCII characters is insufficient to represent the full Vietnamese composite character set, which requires approximately 232 additional atomic tokens covering all phonologically valid combinations of vowel modification diacritics and tonal markers. The vocabulary is extended by replacing the output projection layer with a new linear layer initialized from pretrained weights for shared characters and from random values for new Vietnamese glyphs, allowing the model to retain its pretrained decoding capability for Latin characters while learning the new Vietnamese distributions through fine-tuning. Class-weighted cross-entropy loss is applied to counteract the severe frequency imbalance between common unaccented characters and rare tone-diacritic combinations that are disproportionately critical for lexical correctness. Fine-tuning is conducted on a 2,000,000-image synthetic recognition corpus generated by a three-branch parallel pipeline covering clean printed, augmented scene text, and mixed-background conditions. At inference time, a phonological validity mask restricts decoder outputs to phonologically possible Vietnamese character sequences, reducing the rate of linguistically impossible predictions under degraded image conditions. Together, these adaptations produce a recognition model that outperforms all evaluated baselines on the Vietnamese test set, including systems specifically designed for Vietnamese such as VietOCR \cite{QuocBinhNguyen2020}.
